Existing IoT intrusion detection research has extensively explored
machine learning and deep learning approaches for identifying malicious
network activities. Traditional machine learning based IDS approaches
have demonstrated strong classification capability; however, most
existing methods primarily focus on attack detection without providing
context-aware prioritization of security events.

Alsaedi et al. introduced the TON\_IoT telemetry dataset, enabling
evaluation of machine learning approaches for IoT and industrial IoT
security scenarios \cite{alsaeedi2020toniot}. Several studies have
investigated machine learning and deep learning based intrusion
detection models. Shone et al. proposed a deep learning approach for
network intrusion detection using representation learning techniques
\cite{shone2018deep}. Similarly, Vinayakumar et al. demonstrated the
effectiveness of deep neural architectures for large-scale intrusion
detection problems \cite{vinayakumar2019deep}.

Ensemble learning and anomaly detection approaches have also been
widely explored. Mirsky et al. proposed Kitsune, an ensemble of
autoencoders designed for online network intrusion detection
\cite{mirsky2018kitsune}. Although these approaches improve detection
performance, they generally do not incorporate security risk estimation
or automated response prioritization.

Recent research has explored explainable AI and intelligent security
analytics to improve operational usefulness. Ahmad et al. investigated
machine learning based intrusion detection with explainability
mechanisms \cite{ahmad2021explainable}. However, existing approaches
still lack an integrated framework combining detection, risk scoring,
severity estimation, and decision support.

\begin{table}[H]
\caption{Comparison of RiskAdaptive with Existing Approaches}
\centering
\begin{tabular}{p{2cm}p{1.3cm}p{1.3cm}p{1.3cm}p{1.3cm}}
\toprule
\textbf{Approach} & 
\textbf{Detection} &
\textbf{Risk Score} &
\textbf{Severity} &
\textbf{Response Support}\\
\midrule
ML based IDS \cite{ahmed2016survey} & Yes & No & Limited & No\\
TON\_IoT ML Models \cite{alsaeedi2020toniot} & Yes & No & No & No\\
Deep Learning IDS \cite{vinayakumar2019deep} & Yes & No & No & No\\
Kitsune \cite{mirsky2018kitsune} & Yes & No & No & No\\
RiskAdaptive & Yes & Yes & Yes & Yes\\
\bottomrule
\end{tabular}
\end{table}


Unlike existing IDS approaches, RiskAdaptive introduces an additional
risk-aware decision layer that converts detection outputs into
actionable security priorities. The proposed framework therefore
bridges the gap between intrusion detection and autonomous security
operations by combining threat identification, risk estimation, and
severity-based prioritization.
